% TOP_PROBLEM: {"id": 2303010, "title": "Research Problems in Function Theory — Problem 3.10", "category_id": 8, "category": {"id": 8, "name": "analysis", "display_name": "Analysis", "description": "Limits, continuity, calculus, and function theory.", "slug": "analysis", "order_index": 8, "created_at": "2026-07-31T15:26:25.671Z"}, "status": "open"}
% FIRST_POSTED: null
% ENTRY_KIND: statement_only
% Imported from: batch12.tex; UnsolvedMath ID 2303010
% Compile twice: lualatex 2303010.tex
\documentclass[11pt]{article}
\usepackage{fontspec}
\setmainfont{Latin Modern Roman}
\usepackage{amsmath,amssymb,mathtools,unicode-math}
\setmathfont{Latin Modern Math}
\usepackage{xurl,hyperref}
\hypersetup{colorlinks=true,urlcolor=blue,pdftitle={UnsolvedMath Problem 2303010}}
\newcommand{\sref}[2]{\hyperlink{source:#1:#2}{[S#2]}}
\newcommand{\cataloguescope}[1]{\paragraph{Mathematical question.}#1}
\begin{document}
% UnsolvedMath ID 2303010; source catalogue code AMR-022-3010
\setcounter{section}{2303009}
\section{Research Problems in Function Theory — Problem 3.10}\label{2303010}
{\small\sffamily UnsolvedMath ID 2303010 · Analysis}
\subsection{Definitions and mathematical statement}

\paragraph{Shared notation.}$\mathbb N=\{1,2,\ldots\}$, and $\mathbb Z,\mathbb Q,\mathbb R,\mathbb C$ denote the integers, rationals, reals and complex numbers. Any other conventions are specified locally.


A subharmonic function is an upper semicontinuous function with values in $[-\infty,\infty)$, not identically $-\infty$ on a connected component, that satisfies the mean-value inequality on every closed ball in its domain. A real harmonic function is a twice continuously differentiable function $u$ satisfying $\Delta u=0$. Let $B=\{x\in\mathbb R^3:|x|<1\}$. The condition $u\in C^\infty(\overline B)$ means that derivatives of every order extend continuously to the closed ball. At $x\in\partial B$, the outward normal derivative is $\partial_\nu u(x)=\nabla u(x)\cdot x$.
The source update reports that Wolff constructed counterexamples to this historical question.
\paragraph{Statement recorded in the supplied source.}
If a real harmonic $u\in C^\infty(\overline B)$ has both $u=0$ and $\partial_\nu u=0$ on a subset $E\subset\partial B$ of positive surface area, must $u\equiv0$? \sref{2303010}{2}
\subsection{Short English statement}
Do zero boundary values and zero normal derivatives on a positive-area part of the three-dimensional sphere force a smooth harmonic function to vanish?
\subsection{Sources}
\begin{description}
\item[\hypertarget{source:2303010:1}{S1}] ULAM.AI and the UnsolvedMath Contributors, UnsolvedMath: A Curated Collection of Open Mathematics Problems, record 2303010 (AMR-022-3010). CC BY 4.0. \url{https://huggingface.co/datasets/ulamai/UnsolvedMath}.
\item[\hypertarget{source:2303010:2}{S2}] W. K. Hayman and E. F. Lingham, \emph{Research Problems in Function Theory}, arXiv:1809.07200v2 (2018), Problem 3.10 and Update 3.10. \url{https://arxiv.org/abs/1809.07200}.
\end{description}
\label{end:2303010}
\end{document}
